\documentclass[11pt]{article}
\usepackage{booktabs}
\usepackage{graphicx}
\usepackage{amsmath}
\usepackage{array}
\usepackage[margin=1in]{geometry}
\usepackage{hyperref}

\title{A Software-Defined Network Intrusion Detection System:\\
Deployment-Aligned Random Forest Classification on InSDN and\\
Live ML Detection in Mininet/Ryu}
\author{[Your Name] \\ [Course Code] --- [Institution]}
\date{}

\begin{document}
\maketitle

\begin{abstract}
We present a software-defined network (SDN) intrusion detection system in which a
single machine-learning classifier is both trained offline on the InSDN dataset and
deployed as the sole live detector inside a Ryu controller monitoring a multi-switch
Mininet testbed. The classifier reaches 98.6\% accuracy and 0.988 weighted-F1 on seven
traffic classes; macro-F1 is 0.681, and the gap between the two averages is the report's
central quantitative finding. It is driven by three minority classes: Web-Attack and BFA
sit near F1 0.38 and Botnet at 0.58, because seven-hundred-odd flow entries cannot be
separated from benign traffic on volumetric features alone. The system's design
contribution is the alignment of training and inference. Rather than training on raw
CICFlowMeter micro-flows, the pipeline re-aggregates the dataset into the unidirectional
flow-table entries a controller actually polls, encodes protocols as raw IANA numbers on
both sides, drops byte-derived features that leak the label on InSDN, and adds four
destination-context (fan-in / fan-out) features so that spoofed-source floods --- which
are invisible per flow --- become learnable. The saved artifact is a
\texttt{Pipeline(StandardScaler, RandomForest)}, so no preprocessing is lost at
inference. On the live testbed the model raised 120 alerts against a single hping3 run,
correctly identifying the attacker and victim. We also state the boundaries plainly:
single-source floods faster than InSDN's benign traffic ceiling read as normal by
construction, U2R was too rare to train, and control-plane saturation is undetectable
because it prevents flow entries from ever forming.
\end{abstract}

\section{Introduction and Problem Statement}
Software-defined networking separates the control plane from the data plane, giving a
central controller a global view of network state. That same centralization introduces
an attack surface absent from traditional networks: the controller itself. An intrusion
detection system for SDN must therefore answer two questions that are usually treated
separately --- whether a classifier can distinguish attack traffic from benign traffic,
and whether the features that classifier needs are actually obtainable, and mean the same
thing, when computed from the controller at runtime.

Most reported SDN-IDS work answers only the first. A classifier is trained on a public
flow dataset, high accuracy is reported, and the deployment question is left implicit.
This report treats both, because an earlier iteration of this very system exhibited the
classic failure: a Random Forest that scored well on held-out InSDN rows and labelled
every live flow \texttt{Normal}. The features had the same names on both sides and
different distributions. The core of the work described here is the set of changes that
closed that gap, so that one trained model serves as the live detector rather than a
second, hand-written detector standing in for it.

\paragraph{Research questions.}
\begin{enumerate}
    \item \textbf{RQ1:} How well does a Random Forest trained on InSDN detect each
    traffic class once class imbalance is measured rather than averaged away?
    \item \textbf{RQ2:} What preprocessing is required for a model trained on an offline
    flow dataset to remain usable on live controller observations --- and which of those
    steps actually move the result?
    \item \textbf{RQ3:} Does the trained model transfer to a live Ryu/Mininet deployment,
    and what is detectable in practice versus undetectable by construction?
\end{enumerate}

\section{System Architecture}
The system is a single pipeline of seven source components. Detection is performed by the
machine-learning model alone: the controller computes a feature vector per flow entry,
scores it with the model, and raises an alert when the winning class is not
\texttt{Normal} and its probability clears a confidence threshold. There are no
hand-written threshold rules on the live path.

\begin{table}[h]
\centering
\caption{System components and responsibilities.}
\label{tab:components}
\small
\begin{tabular}{p{3.2cm}p{3.1cm}p{7.3cm}}
\toprule
Component & File & Responsibility \\
\midrule
Emulated network & \texttt{mininet\_topo/} \texttt{topology.py}, \texttt{traffic.py} &
Parametrized 3-tier tree (default: 1 core, 2 aggregation, 4 edge switches; 12 hosts;
role-based IPs) on heterogeneous \texttt{TCLink}s; optional benign background traffic \\
Controller & \texttt{controller/} \texttt{ids\_controller.py} & Ryu application: L2
learning switch with IP-aware flow installation, flow-stats polling every 3\,s, batched
model inference, alert dispatch \\
Feature extraction & \texttt{controller/} \texttt{feature\_extractor.py} & Converts one
\texttt{OFPFlowStatsReply} into feature vectors using the shared contract, including
per-batch fan-in / fan-out counts \\
Shared contract & \texttt{model/} \texttt{feature\_config.py} & Single definition of the
feature vector, imported by both training and the live extractor \\
Model + training & \texttt{model/} \texttt{train\_model.py} & Flow-entry emulation,
leakage checks, Random Forest pipeline, deployment parity check \\
Alert store & \texttt{backend/main.py}, \texttt{backend/database.py} & FastAPI service
over SQLite; records each alert with the model's class and confidence \\
Visualisation & \texttt{dashboard/} \texttt{streamlit\_app.py}, \texttt{topology\_view.py}
& Polls the backend, draws the live topology and lights up attacker and victim \\
\bottomrule
\end{tabular}
\end{table}

\subsection{The shared feature contract}
\label{sec:features}
A single module, \texttt{model/feature\_config.py}, defines the feature vector through one
function, \texttt{build\_feature\_vector()}, which is called identically by the training
script and by the live extractor. The vector's shape, order, and meaning therefore cannot
drift between training and inference by construction; a unit test asserts the live vector
is byte-for-byte the training vector for the same counters.

The contract defines eleven candidate features in a fixed order
(Table~\ref{tab:featurelist}), of which the deployed model uses eight. Three byte-derived
features are dropped at training time (Section~\ref{sec:preproc}). Host- and
application-level features common in older datasets (failed-login counts, TCP-flag
statistics) are deliberately excluded: a Ryu controller cannot compute them from flow
statistics, so a model depending on them would be untrainable in this deployment.

\begin{table}[h]
\centering
\caption{The feature contract. The deployed model uses the eight features not marked
``dropped''; byte-derived features are removed because InSDN does not record byte counts
for attack traffic (Section~\ref{sec:preproc}).}
\label{tab:featurelist}
\begin{tabular}{llp{6.6cm}}
\toprule
\# & Feature & Definition \\
\midrule
1 & \texttt{duration\_sec} & Flow-entry lifetime in seconds (floored at 0.001) \\
2 & \texttt{packet\_count} & Packets matched by the flow entry \\
3 & \texttt{byte\_count} & Bytes matched by the flow entry \emph{(dropped)} \\
4 & \texttt{avg\_packet\_size} & \texttt{byte\_count} / \texttt{packet\_count}
\emph{(dropped)} \\
5 & \texttt{packet\_rate} & \texttt{packet\_count} / \texttt{duration\_sec} (pps) \\
6 & \texttt{byte\_rate} & \texttt{byte\_count} / \texttt{duration\_sec} \emph{(dropped)} \\
7 & \texttt{protocol} & Raw IANA number: 1\,=\,ICMP, 6\,=\,TCP, 17\,=\,UDP, 0\,=\,other \\
8 & \texttt{flows\_to\_dst} & Flow entries in this reply aimed at the same destination \\
9 & \texttt{distinct\_srcs\_to\_dst} & Distinct sources aimed at the same destination
(fan-in) \\
10 & \texttt{flows\_from\_src} & Flow entries in this reply out of the same source \\
11 & \texttt{distinct\_dsts\_from\_src} & Distinct destinations from that source (fan-out) \\
\bottomrule
\end{tabular}
\end{table}

The protocol encoding is a deliberate design point. Both OpenFlow's \texttt{ip\_proto}
match field and InSDN's \texttt{Protocol} column are already IANA numbers, so both sides
use them as-is. An earlier build compacted them to $0/1/2$ on the controller side only,
so every TCP flow reached the forest tagged as protocol $0$ (``unknown'') and was routed
down the wrong branch at every protocol split. A unit test now pins the IANA encoding.

The four destination-context features (8--11) are the design's answer to spoofed-source
floods. InSDN's DDoS traffic forges a distinct source per packet, so each forged source
produces its own flow entry carrying about two packets; on the seven per-flow features
alone, the median DDoS entry is statistically indistinguishable from an idle benign flow,
because it genuinely is one. The attack exists only in the aggregate: one destination
fielding flows from thousands of sources. The fan-in and fan-out counts, computed across
one \texttt{FlowStatsReply}, are exactly the scope the controller has at runtime and make
that aggregate learnable. They remain model features, not thresholds: nothing here decides
``malicious'' on its own.

\section{Dataset and Preprocessing}

\subsection{Source}
InSDN~\cite{insdn} comprises three CSV exports --- normal traffic
(\texttt{Normal\_data.csv}), Metasploitable-target attack traffic
(\texttt{metasploitable-2.csv}), and OVS-collected traffic (\texttt{OVS.csv}) --- each
with 84 CICFlowMeter columns. Concatenated, they total 343,889 flow records.

\subsection{Label consolidation}
Raw labels are normalised to lowercase and mapped to classes by substring matching, with
\texttt{ddos} tested before \texttt{dos} since the latter is a substring of the former.
Two data-quality artifacts surface and are worth recording: the raw \texttt{Label} column
carries \texttt{DDoS} under two distinct strings (73,529 and 48,413 records), and Botnet
is written \texttt{BOTNET}. After consolidation no record falls through to
\texttt{Unknown}. The raw class distribution (Table~\ref{tab:dist}) spans an imbalance
ratio of 7,173:1 between the largest and smallest class.

\begin{table}[h]
\centering
\caption{Raw class distribution after label consolidation, before flow-entry emulation.}
\label{tab:dist}
\begin{tabular}{lrr}
\toprule
Class & Records & Share \\
\midrule
DDoS & 121{,}942 & 35.5\% \\
Probe & 98{,}129 & 28.5\% \\
Normal & 68{,}424 & 19.9\% \\
DoS & 53{,}616 & 15.6\% \\
BFA & 1{,}405 & 0.41\% \\
Web-Attack & 192 & 0.06\% \\
Botnet & 164 & 0.05\% \\
U2R & 17 & 0.005\% \\
\midrule
Total & 343{,}889 & 100\% \\
\bottomrule
\end{tabular}
\end{table}

\subsection{Flow-entry emulation: aligning training with inference}
\label{sec:emulation}
The pipeline does not train on raw CICFlowMeter rows. CICFlowMeter cuts a packet capture
into short bidirectional micro-flows, so a flood appears as tens of thousands of two-packet
rows lasting milliseconds. An OpenFlow flow-table entry is the opposite: one
unidirectional entry per (\texttt{ipv4\_src}, \texttt{ipv4\_dst}, \texttt{ip\_proto}) that
accumulates every matching packet for its whole lifetime, so the same flood is one entry
of tens of thousands of packets over seconds. Same column names, opposite distributions.

\texttt{emulate\_flow\_table\_entries()} re-aggregates the dataset into controller-shaped
entries before any training. It attributes each flow's counters to its
source--destination key, tiles time into windows equal to the controller's flow-entry
lifetime, and emits several samples per entry at evenly spaced poll times, each carrying
the counters accumulated so far and a duration equal to the entry's age at that poll. This
reproduces the 0--10\,s duration range and the cumulative counters the controller reports,
and it doubles as honest augmentation: the model sees each attack both early and fully
developed. Two dataset-specific complications are handled explicitly:

\begin{itemize}
    \item \textbf{Unreliable direction fields.} \texttt{Tot Fwd Pkts} is zero for 100\% of
    DDoS rows and 56\% of rows overall, with every packet booked as ``backward'' even
    though the row is keyed source-to-destination. Splitting on those fields would delete
    the forward half of every flood and turn the attack's fan-in into a fan-out, making
    the victim look like the attacker. When the fields look broken, the total counters are
    attributed to the source-to-destination key instead.
    \item \textbf{Truncated timestamps.} 75\% of timestamps are minute-truncated (the
    seconds field is absent), so rows sharing a truncated timestamp are spread uniformly
    across the interval it represents, in capture order. The spreading is class-agnostic,
    so it does not leak the label; the limitation is up to 60\,s of duration error, which
    makes results sensitive to the chosen flow-timeout.
\end{itemize}

Emulation turns 343,889 records into 301,603 flow-table entries.

\subsection{Byte-feature leakage and downsampling}
\label{sec:preproc}
A byte-counter leakage check runs on the emulated entries. On InSDN the byte count is zero
for 99.6\% of DDoS entries, 72.7\% of Probe entries, and roughly half of BFA, Botnet and
Web-Attack entries, against 23.7\% for Normal. ``The exporter wrote no byte count'' is
therefore very nearly a label, and any byte-derived feature would let the model score well
offline while collapsing on a live testbed where a real switch does count bytes. The three
byte-derived features are dropped, leaving the eight features of Table~\ref{tab:featurelist}.

Two further steps handle the raw imbalance. The dominant class after emulation, spoofed
DDoS, produces 243,886 entries (one per forged source) and is downsampled to 60,000. U2R,
with only 51 emulated entries, falls below the 100-entry minimum and is dropped: seven
volumetric features cannot learn it, and it only corrupts the macro averages. The
resulting seven-class training set is summarised in Table~\ref{tab:entrycounts}.

\begin{table}[h]
\centering
\caption{Emulated flow-table entries, before and after downsampling / minimum-count
dropping. U2R is dropped; DDoS is capped at 60{,}000.}
\label{tab:entrycounts}
\begin{tabular}{lrr}
\toprule
Class & Emulated entries & After preparation \\
\midrule
DDoS & 243{,}886 & 60{,}000 \\
Normal & 55{,}550 & 55{,}550 \\
Probe & 1{,}161 & 1{,}161 \\
DoS & 414 & 414 \\
BFA & 267 & 267 \\
Botnet & 165 & 165 \\
Web-Attack & 109 & 109 \\
U2R & 51 & \emph{dropped} \\
\midrule
Total & 301{,}603 & 117{,}666 \\
\bottomrule
\end{tabular}
\end{table}

\section{Model and Training}
The prepared entries are split 75/25 with stratification, giving 88,249 training and
29,417 test entries. The classifier is a \texttt{Pipeline(StandardScaler,
RandomForestClassifier)} with 150 trees, \texttt{max\_depth=20},
\texttt{class\_weight="balanced"}, and a fixed seed. Saving a pipeline rather than a bare
estimator is deliberate: it guarantees that whatever preprocessing training applied is
carried to inference inside the same artifact. The scaler is a no-op for a scale-invariant
forest, but its presence means a later swap to an SVM or MLP cannot reintroduce the
``where did the scaler go'' defect.

Random Forest suits these features for two reasons. First, feature magnitudes span several
orders (packet counts in the tens of thousands against a protocol code in a small integer
range); a tree ensemble splits each feature on its own scale. Second, the ensemble yields
per-feature importances (Table~\ref{tab:importance}), so a flagged flow can be explained.
Consistent with the design intent, the two rate/count features and duration carry the most
weight, but the four destination-context features together account for about 36\% of total
importance --- confirming that the fan-in / fan-out signal is doing real work, not padding.

\begin{table}[h]
\centering
\caption{Random Forest feature importances (deployed eight-feature model).}
\label{tab:importance}
\begin{tabular}{lr}
\toprule
Feature & Importance \\
\midrule
\texttt{packet\_rate} & 0.233 \\
\texttt{packet\_count} & 0.206 \\
\texttt{duration\_sec} & 0.131 \\
\texttt{distinct\_srcs\_to\_dst} & 0.109 \\
\texttt{flows\_to\_dst} & 0.105 \\
\texttt{distinct\_dsts\_from\_src} & 0.078 \\
\texttt{flows\_from\_src} & 0.071 \\
\texttt{protocol} & 0.067 \\
\bottomrule
\end{tabular}
\end{table}

\subsection{Deployment parity check}
Before saving, training refuses to ship a model that cannot see an obvious attack. It
builds feature vectors at live counter scale --- exactly what the controller produces from
an hping3 flood --- through the same \texttt{build\_feature\_vector()} and asserts they are
not all classified \texttt{Normal}, and that benign probes are not all flagged. This is a
guard against train/serve skew that a test split cannot catch, because the test split
contains no live-scale flood. Its verdicts are recorded in the model metadata and are the
honest basis for the coverage discussion in Section~\ref{sec:live}.

\section{Experiment 1: Per-Class Performance}
\label{sec:perclass}

\begin{table}[h]
\centering
\caption{Per-class results on the held-out set ($n = 29{,}417$).}
\label{tab:perclass}
\begin{tabular}{lrrrr}
\toprule
Class & Precision & Recall & F1 & Support \\
\midrule
DDoS & 1.00 & 1.00 & 1.00 & 15{,}000 \\
Normal & 0.99 & 0.98 & 0.99 & 13{,}888 \\
Probe & 0.77 & 0.89 & 0.82 & 290 \\
DoS & 0.61 & 0.62 & 0.61 & 104 \\
Botnet & 0.45 & 0.83 & 0.58 & 41 \\
BFA & 0.28 & 0.58 & 0.38 & 67 \\
Web-Attack & 0.32 & 0.48 & 0.38 & 27 \\
\midrule
Accuracy & & & 0.986 & 29{,}417 \\
Macro avg & 0.63 & 0.77 & 0.68 & \\
Weighted avg & 0.99 & 0.99 & 0.99 & \\
\bottomrule
\end{tabular}
\end{table}

\subsection{Why the headline number misleads}
Accuracy 0.986 and weighted-F1 0.988 are dominated by DDoS and Normal, which together hold
98.2\% of the test set. Macro-F1 weights all seven classes equally and returns 0.681. The
0.31 spread between weighted- and macro-F1 measures how much of the system's apparent
competence is concentrated in the two classes it was given tens of thousands of examples
of. Reporting both numbers, and explaining the gap, is the substantive finding of this
section. The high-volume classes are separated cleanly: DDoS reaches F1 1.00 (the
destination-context features resolving the spoofed flood), Normal 0.99, and Probe 0.82.
The three small classes carry the macro average down.

\subsection{The minority classes}
BFA (67 test entries, F1 0.38), Web-Attack (27, 0.38), and Botnet (41, 0.58) are the weak
points. Two causes compound. First, sample count: after emulation and the byte-feature
drop, each is represented by a few hundred entries at most, so \texttt{class\_weight}
pushes the boundary toward them and produces the characteristic high-recall,
low-precision signature (Botnet recall 0.83 at precision 0.45). Second, feature overlap:
a web attack is a small number of TCP flows and a brute-force attempt is steady
moderate-rate TCP, neither of which is volumetrically distinct from ordinary traffic once
payload is unobservable. The confusion matrix (Table~\ref{tab:cm}) shows the leakage is
mutual and mostly into and out of Normal, not a single dominating off-diagonal cell.

\begin{table}[h]
\centering
\caption{Confusion matrix. Rows are true classes, columns predicted.}
\label{tab:cm}
\small
\setlength{\tabcolsep}{4pt}
\begin{tabular}{l|rrrrrrr}
\toprule
True $\backslash$ Pred & BFA & Botnet & DDoS & DoS & Normal & Probe & Web-Att. \\
\midrule
BFA        & 39 & 0  & 0      & 3      & 18      & 4     & 3  \\
Botnet     & 0  & 34 & 0      & 0      & 7       & 0     & 0  \\
DDoS       & 0  & 0  & 14{,}979 & 0    & 21      & 0     & 0  \\
DoS        & 2  & 0  & 0      & 64     & 30      & 8     & 0  \\
Normal     & 90 & 42 & 6      & 35     & 13{,}626 & 66   & 23 \\
Probe      & 1  & 0  & 0      & 3      & 26      & 258   & 2  \\
Web-Attack & 5  & 0  & 0      & 0      & 8       & 1     & 13 \\
\bottomrule
\end{tabular}
\end{table}

\subsection{False-positive rate and the confidence threshold}
Because an IDS that cries wolf gets switched off, the false-positive rate --- benign flows
on which an alert fired --- matters operationally more than accuracy. At the deployed
confidence threshold of 0.60, the overall false-positive rate is 1.2\% at a detection
recall of 0.99. Table~\ref{tab:sweep} sweeps the threshold that
\texttt{ids\_controller.py} applies at runtime, and justifies the chosen value rather than
asserting it: raising it to 0.70 cuts the false-positive rate to 0.4\% for one point of
recall, and 0.60 is the knee of that trade-off.

\begin{table}[h]
\centering
\caption{Confidence-threshold sweep. Rows below the threshold are reported \texttt{Normal},
exactly as the controller does. The deployed value is 0.60.}
\label{tab:sweep}
\begin{tabular}{rrrrr}
\toprule
Threshold & Detection recall & False-positive rate & Precision & Accuracy \\
\midrule
0.00 & 0.9929 & 0.0189 & 0.9833 & 0.9863 \\
0.50 & 0.9920 & 0.0179 & 0.9841 & 0.9865 \\
\textbf{0.60} & \textbf{0.9904} & \textbf{0.0122} & \textbf{0.9891} & \textbf{0.9887} \\
0.70 & 0.9876 & 0.0037 & 0.9967 & 0.9915 \\
0.90 & 0.9826 & 0.0006 & 0.9994 & 0.9904 \\
\bottomrule
\end{tabular}
\end{table}

\subsection{Inference latency}
Model inference is measured the way the controller calls it. A single-flow
\texttt{predict\_proba} averages about 56\,ms, dominated by fixed scikit-learn overhead;
scoring a whole \texttt{FlowStatsReply} in one call amortises that to roughly 0.6\,ms per
flow at a batch of 128. The controller uses the batched path, which is why classification
stays comfortably inside the 3-second poll interval. End-to-end detection latency is
dominated by the poll interval, not by inference.

\section{Experiment 2: What Preprocessing Actually Matters}
RQ2 asks which preprocessing steps are load-bearing. The evidence is qualitative but
decisive, because several steps are the difference between a model that deploys and one
that does not, rather than a few points of F1.

\begin{itemize}
    \item \textbf{Flow-entry emulation is necessary for deployment, not accuracy.} The
    pre-emulation model scored about 0.93 offline and labelled every live flow
    \texttt{Normal}, because its decision structure lived entirely in the sub-second
    duration regime that a controller never produces. Emulation (Section~\ref{sec:emulation})
    is what makes the offline and live distributions the same object.
    \item \textbf{Destination-context features are necessary for spoofed DDoS.} Without
    them the model cannot represent a spoofed flood at all, since the signal lives in the
    fan-in / fan-out counts and not in any single flow. With them, DDoS reaches F1 1.00.
    An ablation flag (\texttt{-{}-no-context-features}) reproduces the failing configuration
    for comparison.
    \item \textbf{Dropping byte features is necessary for honesty.} Keeping them lifts
    offline scores toward 1.00 by leaking the capture artifact of a missing byte count,
    and produces a model that cannot see a live flood. The pipeline drops them
    automatically when the leakage check trips.
    \item \textbf{The IANA protocol encoding and the pipeline-saved scaler} are each
    necessary correctness fixes; getting either wrong silently routes every flow to the
    wrong answer at inference time.
\end{itemize}

In short, the minority-class weakness in Section~\ref{sec:perclass} is a genuine
data-and-feature-space limit, whereas the deployment failure of the earlier model was a
train/serve alignment problem that preprocessing solved outright.

\section{The Live Detection Path}
\label{sec:live}
The trained model is the live detector. The controller installs one unidirectional flow
entry per (source, destination, protocol) with a 10-second hard timeout --- matching the
flow-timeout the model was trained on --- polls flow statistics every 3\,s, builds the
feature vectors for the whole reply through the shared extractor, scores them in one
batched call, and raises an alert when the winning class is not \texttt{Normal} and its
probability clears 0.60. A per-victim cooldown suppresses repeat alerts so one sustained
attack does not bury the feed. The controller performs no mitigation: it classifies and
reports.

A legacy threshold detector (\texttt{controller/detector.py}) remains in the repository
from the earlier build but is not imported by the controller and plays no part in live
detection; it is noted here only so its presence is not mistaken for a second active path.

\subsection{Live testbed results}
Against a single hping3 SYN-flood run from attacker host \texttt{h6} (10.1.2.3) to the web
server \texttt{h1} (10.1.1.1), the model raised 120 alerts over about seven minutes,
recorded in the backend database. The breakdown is 86 DoS and 34 Probe alerts, with
confidences from 0.62 to 0.95, and a mean live inference time of 8.5\,ms per flow. The
alerts correctly name the attacker as source and the victim as destination, and appear
across the edge and aggregation switches on the path. The average packet size in these
alerts is 54--58 bytes, consistent with the near-empty packets of a SYN flood. A handful
of early alerts carry MAC addresses in place of IPs, from flows seen before the switch had
learned the hosts; these are cosmetic and clear once L2 learning completes.

\subsection{What is detectable, and what is not}
The offline deployment-parity check and the live run together bound what the system can
see. The model detects spoofed DDoS (via the context features), normal traffic (correctly
silent), and floods whose rate falls within the range InSDN's DoS traffic covers --- which
is where the Mininet testbed's link-limited hping3 floods landed, hence the 120 live DoS
and Probe alerts. It does \emph{not} detect single-source floods faster than InSDN's
benign ceiling: the parity check shows a 30{,}000-packet, 3-second SYN flood classified
\texttt{Normal} at confidence 0.59, because InSDN's benign bulk transfers reach 311{,}000
pps while its DoS traffic tops out around 6{,}555 pps, so at the highest flood rates the
training data itself says ``Normal''. This is a dataset limit, recorded in the model
metadata as an uncovered shape, and it is honest to report it rather than hide it behind
the aggregate accuracy.

\section{Attack Taxonomy and Coverage}
The simulated attacks, their tools, the InSDN classes they correspond to, and the live
model output are reconciled in Table~\ref{tab:taxonomy}.

\begin{table}[h]
\centering
\caption{Attack coverage across the testbed.}
\label{tab:taxonomy}
\small
\begin{tabular}{p{2.9cm}p{2.6cm}p{2.3cm}p{5.0cm}}
\toprule
Simulated attack & Tool & InSDN class & Live model output \\
\midrule
SYN flood & \texttt{hping3 -S --flood} & DoS / DDoS & DoS (observed live), or DDoS when the
context features see multiple sources \\
UDP flood & \texttt{hping3 --udp --flood} & DoS / DDoS & DoS / DDoS by the same signal \\
Port scan & \texttt{nmap -sS -p-} & Probe & Probe (observed live) \\
Slowloris & \texttt{slowloris\_attack.py} & \emph{no clean analogue} & Low-and-slow is
weakly represented in the volumetric features; not reliably distinguished \\
Control-plane saturation & \texttt{control\_plane\_saturation.py} & \emph{absent from
InSDN} & \emph{not detected by construction} \\
\bottomrule
\end{tabular}
\end{table}

Control-plane saturation is outside the detection path by construction. The attack
randomises source IP and MAC on every packet so that no flow-table entry ever matches,
forcing a \texttt{PacketIn} per packet and targeting the controller's processing capacity
rather than any host. The model has no training examples of it, because InSDN contains no
such class, and the live detector cannot see it either, because detection operates on
flow-statistics polling and the attack is designed to prevent flow entries from forming.
Detecting it would require a \texttt{PacketIn} arrival-rate counter in the controller,
independent of the flow-stats path --- instrumentation the system does not have.

\section{Discussion}
The result that generalises beyond this dataset is the alignment argument. A classifier's
offline accuracy is evidence about deployability only when the features it consumes mean
the same thing at inference time. Here that required re-deriving the training features from
the shape the controller actually observes --- unidirectional, cumulative, IANA-encoded,
byte-blind, and context-aware --- rather than from the exporter's native micro-flows. Once
aligned, one model trained on a public dataset served as the live detector and produced
correct attacker/victim attributions on a testbed it was never trained on.

The minority-class weakness is the honest cost. Web-Attack, BFA, and Botnet cannot be
separated from benign traffic on volumetric flow features, and no amount of reweighting
recovers a signal that lives in packet payload the feature space never observes. The
appropriate response is to scope the system's claims to the classes it can support ---
high-volume floods, spoofed DDoS, and port scans --- and to report the minority F1 scores
as the measured limits they are.

\section{Limitations}
\begin{enumerate}
    \item \textbf{Minority classes are weak.} Web-Attack and BFA sit near F1 0.38 and
    Botnet at 0.58, from a few hundred entries each and volumetric feature overlap
    (Section~\ref{sec:perclass}).
    \item \textbf{U2R is not modelled.} 17 raw records became 51 emulated entries, below
    the training minimum, so the class was dropped rather than reported at a noise-level
    metric.
    \item \textbf{High-rate single-source floods above InSDN's benign ceiling read as
    Normal.} A documented dataset gap, recorded in the model metadata; closing it requires
    training on captured Mininet flow statistics.
    \item \textbf{Control-plane saturation is undetectable by construction}, as it prevents
    flow entries from forming and no independent \texttt{PacketIn} counter exists.
    \item \textbf{Byte-derived features are unavailable} on InSDN because the dataset does
    not record byte counts for attack traffic; the model runs on eight non-byte features.
    \item \textbf{Timestamp reconstruction} introduces up to 60\,s of duration error, so
    results are sensitive to the flow-timeout and should be reported with that sensitivity.
    \item \textbf{Thresholds and topology are calibrated to this testbed} (a 12-host,
    7-switch tree with mixed link profiles) and carry no claim of transfer to other scales
    or link speeds.
    \item \textbf{Fat-tree mode needs a loop-aware controller.} The current L2 learning
    switch floods, so redundant paths would broadcast-storm; only the loop-free tree mode
    is used for results.
\end{enumerate}

\section{Conclusion}
A single Random Forest, trained on InSDN and deployed unchanged inside a Ryu controller,
reaches 98.6\% accuracy and 0.988 weighted-F1 but 0.681 macro-F1; all three numbers belong
in the same sentence, because the first two are carried by two large classes and the third
records that three of seven are hard. The system's contribution is not the headline score
but the alignment that makes it deployable: re-aggregating the dataset into
controller-shaped flow entries, encoding protocols consistently, dropping features that
leak a capture artifact, and adding destination-context features that turn an
otherwise-invisible spoofed flood into a class with F1 1.00. On the live testbed the model
raised 120 correctly attributed alerts against a real attack. Its boundaries are stated
precisely: the rarest classes, the fastest single-source floods, and the attack most
specific to SDN --- control-plane saturation --- lie outside what flow statistics can
show. Stating those boundaries is what distinguishes a measured system from a
demonstration.

\begin{thebibliography}{9}
\bibitem{insdn} Elsayed, M. S., Le-Khac, N.-A., Dev, S., \& Jurcut, A. D. (2020). InSDN:
A Novel SDN Intrusion Dataset. \emph{IEEE Access}, 8, 165263--165284.
\bibitem{rf} Breiman, L. (2001). Random Forests. \emph{Machine Learning}, 45(1), 5--32.
\bibitem{ryu} Ryu SDN Framework Community. \emph{Ryu: Component-Based Software Defined
Networking Framework}. \url{https://ryu-sdn.org}
\bibitem{mininet} Lantz, B., Heller, B., \& McKeown, N. (2010). A Network in a Laptop:
Rapid Prototyping for Software-Defined Networks. \emph{Proc. 9th ACM SIGCOMM Workshop on
Hot Topics in Networks (HotNets)}.
\bibitem{sklearn} Pedregosa, F. et al. (2011). Scikit-learn: Machine Learning in Python.
\emph{Journal of Machine Learning Research}, 12, 2825--2830.
\end{thebibliography}

\end{document}
